\documentclass[11pt]{article}
\usepackage[utf8]{inputenc}
\usepackage[spanish, es-tabla]{babel}
\usepackage{amsmath, amssymb}
\usepackage{enumitem}
\usepackage{geometry}
\usepackage{graphicx}
\geometry{letterpaper, margin=2.5cm}

\title{Econometría ICS-294\\Ejercicios: Repaso Estadística II}
\author{Francisco Alfaro Medina}
\date{}

\begin{document}
\maketitle

\section*{Instrucciones}
Responda de forma ordenada, justificando cada fórmula utilizada. Cuando corresponda, interprete el resultado en palabras.

\section*{Ejercicios}

\begin{enumerate}[label=\textbf{Ejercicio \arabic*.}, leftmargin=*]
    \item \textbf{Teorema Central del Límite.}
    
    Una variable aleatoria de interés tiene media poblacional $\mu=50$ y desviación estándar $\sigma=12$. Se toman muestras aleatorias de tamaño $n=36$.
    \begin{enumerate}[label=\alph*)]
        \item ¿Cuál es la media de la distribución muestral de $\bar{X}$?
        \item ¿Cuál es la desviación estándar de la distribución muestral de $\bar{X}$?
        \item Explique por qué podemos aproximar la distribución de $\bar{X}$ usando una normal.
    \end{enumerate}

    \item \textbf{Intervalo de confianza para una media con $\sigma$ conocida.}
    
    Una empresa quiere estimar el gasto mensual promedio de sus clientes. En una muestra aleatoria de $n=64$ clientes se obtiene un gasto promedio de $\bar{x}=85{,}000$ pesos. Suponga que la desviación estándar poblacional es conocida e igual a $\sigma=20{,}000$.
    \begin{enumerate}[label=\alph*)]
        \item Construya un intervalo de confianza de 95\% para $\mu$.
        \item Construya un intervalo de confianza de 99\% para $\mu$.
        \item Compare ambos intervalos. ¿Cuál es más preciso? ¿Cuál es más confiable?
    \end{enumerate}

    \item \textbf{Interpretación de un intervalo de confianza.}
    
    Un investigador obtiene el siguiente intervalo de confianza de 95\% para el ingreso promedio mensual de una población:
    \[
        IC_{95\%}=(720{,}000;\, 780{,}000).
    \]
    \begin{enumerate}[label=\alph*)]
        \item Interprete correctamente este intervalo.
        \item Explique por qué no es correcto decir: ``hay 95\% de probabilidad de que $\mu$ esté en este intervalo'', una vez calculado el intervalo.
    \end{enumerate}

    \item \textbf{Error absoluto y precisión.}
    
    Suponga que se quiere estimar una media poblacional con $\sigma=15$, usando un intervalo de confianza de 95\%.
    \begin{enumerate}[label=\alph*)]
        \item Si $n=100$, calcule el error absoluto $E$.
        \item Si se quiere que el error absoluto no supere $E=2$, ¿qué tamaño muestral mínimo se requiere?
        \item ¿Qué ocurre con $n$ si se exige un intervalo de 99\% manteniendo el mismo error máximo?
    \end{enumerate}

    \item \textbf{Determinación del tamaño muestral.}
    
    Una municipalidad quiere estimar el tiempo promedio de viaje al trabajo de sus habitantes. Estudios previos sugieren que $\sigma=18$ minutos. Se desea un intervalo de confianza de 95\% con error máximo de 3 minutos.
    \begin{enumerate}[label=\alph*)]
        \item Calcule el tamaño muestral requerido.
        \item Si se quisiera un error máximo de 2 minutos, ¿cómo cambia el tamaño muestral?
        \item Explique la relación entre precisión y tamaño muestral.
    \end{enumerate}

    \item \textbf{Intervalo de confianza para una proporción.}
    
    En una encuesta a $n=400$ estudiantes, 256 declaran usar transporte público para llegar a la universidad.
    \begin{enumerate}[label=\alph*)]
        \item Calcule $\hat{p}$.
        \item Verifique si se cumplen las condiciones para usar la aproximación normal.
        \item Construya un intervalo de confianza de 95\% para la proporción poblacional.
        \item Interprete el intervalo.
    \end{enumerate}

    \item \textbf{Evaluación de una afirmación usando un IC para proporciones.}
    
    Una empresa afirma que solo 8\% de sus productos presentan fallas. Una muestra aleatoria de 300 productos encuentra 36 defectuosos.
    \begin{enumerate}[label=\alph*)]
        \item Construya un intervalo de confianza de 95\% para la proporción de productos defectuosos.
        \item ¿El valor afirmado por la empresa cae dentro del intervalo?
        \item ¿Qué concluye respecto de la afirmación de la empresa?
    \end{enumerate}

    \item \textbf{Tamaño muestral para una proporción.}
    
    Se quiere estimar la proporción de personas que apoyan una política pública. No existe información previa sobre la proporción poblacional.
    \begin{enumerate}[label=\alph*)]
        \item Use $p^*=0.5$ para calcular el tamaño muestral necesario si se desea un intervalo de 95\% con error máximo $E=0.04$.
        \item Repita el cálculo para un intervalo de 99\%.
        \item Explique por qué se usa $p^*=0.5$ cuando no hay información previa.
    \end{enumerate}

    \item \textbf{Diferencia de medias con varianzas conocidas.}
    
    Se quiere comparar el salario promedio de egresados de dos carreras. Para la carrera A se observa $\bar{x}=950{,}000$, $n_x=100$ y $\sigma_x=80{,}000$. Para la carrera B se observa $\bar{y}=890{,}000$, $n_y=121$ y $\sigma_y=70{,}000$.
    \begin{enumerate}[label=\alph*)]
        \item Construya un intervalo de confianza de 95\% para $\mu_x-\mu_y$.
        \item ¿El intervalo contiene al cero?
        \item Interprete el resultado.
    \end{enumerate}

    \item \textbf{Discusión conceptual.}
    
    Responda brevemente:
    \begin{enumerate}[label=\alph*)]
        \item ¿Por qué al aumentar el nivel de confianza aumenta el largo del intervalo?
        \item ¿Por qué aumentar el tamaño muestral mejora la precisión?
        \item ¿Qué significa que dos muestras sean independientes al construir un intervalo para diferencia de medias?
        \item ¿Qué diferencia hay entre estimador puntual e intervalo de confianza?
    \end{enumerate}
\end{enumerate}

\end{document}
